\documentclass[authoryear,1p,12pt]{elsarticle}
\usepackage{amsmath}
\usepackage{amssymb}
\usepackage[mathlines,displaymath]{lineno}
\usepackage{graphicx}
\usepackage[all,2cell,dvips]{xy}
\usepackage{natbib}
\usepackage[usenames,rgb]{xcolor} 
\usepackage{ucs}                                            
\usepackage[utf8]{inputenc}                                                                           
\usepackage{array}                                            
\usepackage{longtable}                                        
\usepackage{calc}                                             
\usepackage{multirow}                                         
\usepackage{hhline}                                           
\usepackage{ifthen} 
% \usepackage{babel}
\usepackage{tcolorbox}
\newtcolorbox{mybox}{colback=grey!5!white,colframe=red!75!black}
\usepackage{setspace}
\usepackage{verbatim}
\usepackage{etaremune}
\usepackage{url}
\usepackage{color}
\usepackage{graphicx}
\usepackage{natbib}
\usepackage[latin1]{inputenc}
\usepackage{color}
\usepackage{array}
\usepackage{float}
\usepackage{wrapfig}

\newcommand{\etal}{{et~al.{}}}
\newcommand{\ie}{{i.~e.{}}}
\newcommand{\eg}{{e.~g.{}}}
\newcommand{\viz}{{viz.{}}}
\newcommand{\etc}{{etc.{}}}
\newcommand{\apriori}{{a priori{}}}
\newcommand{\vv}{{vice versa{}}}
\newcommand{\cf}{cf.{}}
\setcounter{tocdepth}{3}
\usepackage{tikz}
\newcommand{\carlos}[1]{\textcolor{Red}{#1}}

% Text layout
\topmargin 0.0cm
\oddsidemargin 0.2cm
\evensidemargin 0.2cm
\textwidth 17cm 
\textheight 19cm


\begin{document}

%\begin{mybox}\begin{singlespace}
%\begin{small}
\section*{Box 1: BAM-trait based biodiversity model}

We outline a coevolutionary IBM model to evaluate the effect of trait
architecture (i.e., the covariation among biotic, abiotic and
migration traits) on rescue and biodiversity dynamics (Figure 1a). We
explore two trait architectures rooted in empirical and theoretical
studies (Armbruster et al. 2014): the modular and the trait
integration architectures. The modular architecture takes into account
independent traits, i.e. each trait contributes independently to
fitness (Figure 1b: Modular). The integration trait architecture
incorporates the biotic trait as the main contributor to fitness with
the abiotic and the migration traits depending on the evolution of the
biotic trait (Fig 1b:Magic).

Our model considers a landscape with a network of patches connected by
migration events and inhabited by resource and consumer populations
comprising several individuals each. Biotic, abiotic and migration
trait values are initially sampled from normal distributions for each
individual phenotype i represented as {\bf z} = [$z_{b}$, $z_{a}$,
$z_{d}$]. The local resource demography of species $i$ is driven by
the temporal dependent fitness function following
\vspace{0.25 in}

  \begin{equation}
    \mathrm{W_{R}({\bf z_{B_{i}}})}^{t}_{x} = 1 / [\left 1 + exp(\left -\gamma(\mathrm{\bf z_{B_{i}}}^{t}_{x} - {\bf \hat{z}_{B_{for all j}}}^{t}_{x})^2)\right],
  \end{equation}
  \begin{equation}
    \mathrm{W_{R}({\bf z_{A_{i}}})}^{t}_{x} = exp[- \gamma(\mathrm{\bf z_{A_{i}}}^{t}_{x} - \mathrm{\bf \hat{z}_{A_{i}}}^{t}_{x})^2],
  \end{equation}
  \begin{equation}
      \mathrm{W_{R}({\bf z_{M_{i}}})}^{t}_{x} = exp[- \gamma(\mathrm{\bf z_{M_{i}}}^{t}_{x} - \mathrm{\bf \hat{\theta}_{D}}_{x})^2],
    \end{equation}
    \vspace{0.25 in}
    
and for the consumers, the fitness functions are the following

\vspace{0.25 in}
  \begin{equation}
  \mathrm{W_{C}({\bf z_{B_{i}}})}^{t}_{x} = exp[\left -\gamma(\mathrm{\bf z_{B_{i}}}^{t}_{x} - {\bf z_{B_{j}}}^{t}_{x})^2 \right],
  \end{equation}
\begin{equation}
 \mathrm{W_{C}({\bf z_{A_{i}}})}^{t}_{x} = exp[\left - \gamma(\mathrm{\bf z_{A_{i}}}^{t}_{x} - \mathrm{\bf \theta_{A}}_{x})^2 \right],
\end{equation}
\begin{equation}
 \mathrm{W_{C}({\bf z_{M_{i}}})}^{t}_{x} = exp[\left - \gamma(\mathrm{\bf z_{M_{i}}}^{t}_{x} - \mathrm{\bf \theta_{D}}_{x})^2 \right],
  \end{equation}

  \vspace{0.25 in}
where $\mathrm{\bf z_{T}}^{t}_{x}$ is the trait value $T$ of
phenotype $i$ at time $t$ and site $x$, ${\bf \theta}_{x}$ is the
multivariate fitness optimum, ${\bf \omega}$ is the covariance matrix
(Fig 1) \citep{Lande:1980, Melo&Marroig:2014}, and $\gamma$ determines
the interaction sensitivity to deviations from the biotic (i.e.,
coevolutionary selection strength), abiotic and dispersal optimum.

In this equation, demographic and trait evolutionary dynamics depend
on the scenario of trait architecture being considered.

If the covariance matrix, ${\bf \omega}$, is diagonal, representing
our modular scenario (Fig. 1), then there is no correlated stabilizing
selection. If the covariance matrix is not diagonal and considering
our magic scenario (Fig. 1), then there is correlated stabilizing
selection with the biotic trait determining the selection on the other
traits. The biotic optimum changes in time and we define it as the
minimum and maximum mean trait distance across all populations in site
$x$ for the consumer and the resource biotic trait value,
respectively. The dispersal optimum is given by the mean distance in
the landscape (i.e., each site has coordinates from which we calculate
the Euclidian distance between each pair of sites). Phenotypes with
dispersal trait value near to this optimum maximize fitness for this
trait only for the modular but not for the magic trait scenario. The
abiotic optimum is taken initially as the mean of the distribution and
it remains the same during the dynamics.

At each time step a certain percentage of individuals are selected
randomly to die and are replaced by birth or migration. The time
dependent fitness function in site $x$, $\mathrm{W({\bf z})}^{t}_{x}$,
determines the mother of the offspring. To obtain the phenotypic value
for the offspring, {\bf $z_o$} = $[z_{ob}, z_{oa}, z_{od}]$, we add to
each trait value of the mother the environmental deviation, {\bf e},
taken from a multivariate uniform distribution with range $\pm0.001$
in all dimensions (i.e., independent sampled values for the modular
but correlated for the integration trait scenario). We run our model
for 10000 replicates with 50 generations each for the modular and
integration trait scenarios. To evaluate the importance of coevolution
and migration as rescue mechanisms in modulating biodiversity
dynamics, we explore how the formation of hot and cold spots of
species persistence in local and regional biodiversity changes in a
gradient of migration and coevolutionary selection strength (Figure
x2).

\begin{comment}
Mean species richness decays with the strength of coevolutionary
selection for the biotic trait in the two scenarios, the modular
(black) and the integrated (red) scenario (Figure 1b). Adding
interactions among traits, the integrated scenario, decreases the
number of species for all the gradient of coevolutionary selection,
compared to the modular scenario.

Q2: mention the initial values: number of species, population size,
initial trait variance and covariance values, and the number of
generations per replicate.
\end{comment}

Codes and outputs can be downloaded from
https://github.com/melian009/Cometa).
% \end{small}
%\end{singlespace}
%\end{mybox}

%\newpage
\section{Bibliography}
\bibliographystyle{unsrtnat}
\bibliography{space.bib}
\end{document}